\chapter{Introducción}

\section{Motivación y Contexto}

Uno de los desafíos que enfrenta el desarrollo de software moderno es cerrar la brecha entre la especificación del comportamiento deseado de un sistema y la generación automática de casos de prueba que validen ese comportamiento. El Desarrollo Dirigido por Comportamiento (BDD, \textit{Behavior-Driven Development}) ha demostrado ser efectivo para capturar requisitos mediante escenarios en lenguaje natural estructurado, utilizando la sintaxis Given-When-Then del formato Gherkin. Esta aproximación facilita la colaboración entre técnicos y no técnicos, permitiendo que analistas de negocio, testers y desarrolladores compartan un lenguaje común para describir el comportamiento esperado del sistema.

Por otro lado, las Pruebas Basadas en Modelos (MBT, \textit{Model-Based Testing}) ofrecen un enfoque riguroso y automatizado para generar casos de prueba a partir de modelos formales del sistema, como los Sistemas de Transición Simbólicos (STS, \textit{Symbolic Transition Systems}). Estos modelos formales permiten aplicar algoritmos de generación automática de tests de acuerdo con criterios de cobertura explícitos y detectar comportamientos inesperados.

Sin embargo, existe una brecha entre estos dos mundos. Los escenarios BDD, expresados en lenguaje natural, carecen de la precisión formal necesaria para alimentar directamente las herramientas de MBT. Investigaciones recientes han propuesto el Lenguaje Intermedio de BDD (IBDD) como puente entre ambos paradigmas~\cite{zameni2024ibdd}. IBDD mantiene la estructura Given-When-Then de BDD pero incorpora elementos formales de los STS, incluyendo variables globales y locales, guardas lógicas, gates de interacción y asignaciones explícitas de estado.

Los trabajos de Zameni et al. han demostrado que, una vez que los escenarios están representados en IBDD, es posible transformarlos automáticamente en Sistemas de Transición BDD (BDDTS) y posteriormente en STS, desde donde se pueden generar casos de prueba completos y formalmente válidos~\cite{zameni2023bddts}. Sin embargo, un eslabón de esta cadena de automatización aún no ha sido abordado: la traducción inicial de escenarios Gherkin escritos por humanos a representaciones IBDD formalmente correctas. La Figura~\ref{fig:pipeline-overview} ilustra este pipeline y señala la etapa que aborda el presente trabajo.

\vspace{0.5cm}

\begin{figure}[htbp]
\centering
\resizebox{\textwidth}{!}{%
\begin{tikzpicture}[
    node distance=1cm,
    stage/.style={
        rectangle, rounded corners=4pt, draw=black, thick,
        minimum height=1cm, minimum width=2cm,
        text centered, font=\small
    },
    arr/.style={-{Stealth[length=5pt]}, thick},
    lbl/.style={font=\footnotesize\itshape, text=black!70}
]
    \node[stage] (gherkin) {Gherkin};
    \node[stage, right=of gherkin] (ibdd) {IBDD};
    \node[stage, right=of ibdd] (bddts) {BDDTS};
    \node[stage, right=of bddts] (sts) {STS};
    \node[stage, right=of sts, text width=1.8cm, align=center] (tests) {Casos de prueba};

    \draw[arr] (gherkin) -- (ibdd);
    \draw[arr] (ibdd) -- (bddts);
    \draw[arr] (bddts) -- (sts);
    \draw[arr] (sts) -- (tests);

    % Brace for "Este trabajo"
    \draw[decorate, decoration={brace, amplitude=5pt, mirror}, thick]
        ([yshift=-0.6cm]gherkin.south west) -- ([yshift=-0.6cm]ibdd.south east)
        node[midway, below=6pt, lbl] {Este trabajo};

    % Brace for "Formalmente definido"
    \draw[decorate, decoration={brace, amplitude=5pt, mirror}, thick]
        ([yshift=-0.6cm]ibdd.south east) -- ([yshift=-0.6cm]tests.south east)
        node[midway, below=6pt, lbl] {Formalmente definido~\cite{zameni2024ibdd, zameni2025composition}};
\end{tikzpicture}%
}
\caption{Pipeline de transformación BDD-a-MBT y ubicación del presente trabajo}
\label{fig:pipeline-overview}
\end{figure}

Actualmente, esta traducción es un proceso enteramente manual que requiere experiencia especializada en teoría de sistemas de transición y conocimiento profundo de la gramática IBDD. Un analista debe identificar variables implícitas, inferir gates de interacción apropiados, establecer guardas lógicas correctas y determinar las asignaciones de estado necesarias, información que típicamente no está explícita en el texto natural del escenario BDD. Este cuello de botella limita la escalabilidad industrial de todo el pipeline BDD-a-MBT.

\section{Problema}

El problema central que aborda este trabajo puede formularse así:

\begin{quote}
\centering
\textit{¿Cómo automatizar la traducción de escenarios BDD en formato Gherkin al lenguaje intermedio formal IBDD?}
\end{quote}

Esto requiere preservar razonablemente el comportamiento especificado y generar representaciones formalmente válidas para alimentar herramientas de generación automática de casos de prueba. Este problema presenta múltiples desafíos técnicos. En primer lugar, los escenarios Gherkin contienen imprecisiones inherentes al lenguaje natural que deben resolverse para producir especificaciones formales. A esto se suma que elementos cruciales de IBDD (variables globales, tipos de datos, funciones del dominio, gates del sistema) no están explícitos en el texto Gherkin y deben inferirse del contexto o del conocimiento del dominio. La correspondencia entre construcciones lingüísticas en Gherkin y elementos formales de IBDD tampoco es trivial; por ejemplo, la frase ``el usuario hace clic en agregar al carrito'' debe traducirse a un gate de entrada específico con variables de interacción apropiadas. Las traducciones generadas, a su vez, deben ser sintácticamente correctas según la gramática IBDD y semánticamente consistentes con la especificación BDD original. Finalmente, la diversidad de dominios que abarcan los escenarios BDD introduce vocabularios y patrones de interacción específicos que el sistema debe manejar.

\section{Enfoque propuesto}

Este trabajo propone un enfoque híbrido que combina las capacidades de comprensión semántica de Modelos de Lenguaje Grandes (LLMs, \textit{Large Language Models}) con la estructura formal proporcionada por la gramática IBDD y técnicas de análisis sintáctico determinístico.

La arquitectura del sistema se organiza en torno a dos componentes principales. El primero es un módulo de traducción basado en LLM que utiliza modelos de lenguaje para interpretar escenarios Gherkin y generar candidatos de traducción IBDD, aprovechando la capacidad de los LLMs para capturar patrones semánticos complejos y realizar inferencias sobre información implícita. El segundo componente es un módulo de validación sintáctica para verificar las traducciones generadas contra la gramática formal IBDD y garantizar corrección estructural.

El flujo del sistema combina una etapa de traducción de escenarios BDD en formato Gherkin asistida por LLM con una etapa de validación sintáctica mediante el analizador Lark. Cuando una traducción no supera la validación, se activa un ciclo de corrección iterativa orientado a mejorar la conformidad con la gramática IBDD.

La elección de LLMs se fundamenta en su demostrada capacidad para tareas de traducción entre lenguajes formales y naturales, como se evidencia en trabajos recientes sobre generación de código desde especificaciones~\cite{jiang2024codellm} y traducción de requisitos a modelos formales~\cite{soeken2012nlp}. Sin embargo, los LLMs por sí solos no garantizan corrección sintáctica ni consistencia formal, de ahí la necesidad del componente de validación.

\section{Contribuciones}

Este trabajo contribuye con el diseño e implementación de un pipeline automatizado para la traducción BDD$\rightarrow$IBDD que integra generación con LLM, validación sintáctica determinística y corrección iterativa de traducciones fallidas. También formaliza una estrategia de validación operativa basada en la gramática IBDD y parser Earley (Lark), con reportes estructurados de error reutilizables en el ciclo de reintento.

En el plano experimental, se construye y documenta un protocolo reproducible con partición fija del conjunto de datos en desarrollo y prueba, múltiples corridas por configuración y métricas agregadas para controlar la estocasticidad del modelo. Sobre esa base, se realiza una evaluación cuantitativa en cuatro configuraciones lingüísticas de entrada/salida sobre el conjunto de prueba, permitiendo analizar validez inicial y final, ganancia por corrección y desempeño del proceso de corrección. Finalmente, el trabajo delimita el alcance actual del enfoque, analizando calidad sintáctica y señalando la fidelidad semántica como eje de extensión futura.
